\begin{chapintro}
メタ解析の結果は、どの研究を集めるかで決まり、どの研究を集めるかは問いの立て方で決まります。この章では、降圧薬 A の例で問いを PICO の形に書き下し、組み入れ基準と除外基準、解析の計画を、データを見る前に文書にしておく方法を説明します。
\end{chapintro}

\section{問いを PICO に分ける}
\label{sec:02-1}

「降圧薬 A は効くのか」という問いは、そのままでは文献を探すことも選ぶこともできません。誰に対して、何と比べて、何が変わるのかが決まっていないからです。介入の効果を調べる系統的レビューでは、問いを 4 つの要素に分けて書くのが標準のやり方で、それぞれの頭文字をとって PICO と呼びます（図 2.1）。

\begin{description}
\item[P（Population、対象）] どんな人を対象にするのか。降圧薬 A の例では、高血圧のある成人とします。
\item[I（Intervention、介入）] 何を行うのか。例題では降圧薬 A の投与です。
\item[C（Comparison、比較）] 何と比べるのか。例題ではプラセボとします。
\item[O（Outcome、アウトカム）] 何で効果を測るのか。例題では脳卒中の発症とします。
\end{description}

\begin{kkcenter}
\begin{kkfitpic}
\begin{tikzpicture}[
  pb/.style={kbox, minimum width=24mm, minimum height=14mm, align=center, font=\footnotesize},
  node distance=3mm]
\node[pb] (p) {P 対象\\[1mm]高血圧の成人};
\node[pb, right=of p] (i) {I 介入\\[1mm]降圧薬 A};
\node[pb, right=of i] (c) {C 比較\\[1mm]プラセボ};
\node[pb, right=of c] (o) {O アウトカム\\[1mm]脳卒中の発症};
\node[kplain, below=10mm of $(i)!0.5!(c)$, minimum width=80mm, font=\footnotesize, align=center]
  (q) {高血圧の成人で、降圧薬 A はプラセボと比べて\\脳卒中の発症を減らすか};
\foreach \x in {p,i,c,o} \draw[kar] (\x.south) -- (q.north -| \x.south);
\end{tikzpicture}
\end{kkfitpic}
\kkcap{図 2.1　降圧薬 A の問いを PICO の 4 要素に分ける}
\end{kkcenter}

4 つの要素を決めると、問いは「高血圧の成人で、降圧薬 A はプラセボと比べて脳卒中の発症を減らすか」という一文になります。この一文が、文献の検索式（\kkref{sec:03-2}）、組み入れの判断（\kkref{sec:03-3}）、データの抽出項目（\kkref{sec:04-1}）のすべての出発点になります。

要素の決め方によって、集まる試験は大きく変わります。比較を「プラセボ」とすれば、降圧薬 A と別の降圧薬を比べた試験は入りません。比較を「プラセボまたは他の降圧薬」とすれば試験は増えますが、プラセボと比べた効果と、他の降圧薬と比べた効果は意味の違う数字なので、一緒に平均してよいかという問題が生じます。\kkref{sec:01-4} で述べた、違うものを混ぜる問題は、この段階で問いを絞ることで避けます。

\section{アウトカムの決め方}
\label{sec:02-2}

アウトカムは、検査でいえば測定項目にあたります。何を測るのかが曖昧だと、試験ごとに違うものを数えた結果を平均することになります。

アウトカムには、患者にとって意味のある出来事（脳卒中、死亡、入院）と、その代わりに測る指標（血圧の値）があります。血圧が下がることと脳卒中が減ることは関連していますが、同じではありません。血圧をよく下げても脳卒中が減らない薬剤があれば、血圧だけを見たレビューは誤った安心を与えます。そのため、レビューの主要アウトカムには、患者にとって意味のある出来事を選ぶのが原則です。この例では脳卒中の発症を主要アウトカムとし、血圧の変化や有害事象は副次アウトカムとして別に扱います。

アウトカムの定義も先に決めておきます。脳卒中を「画像検査で確認されたもの」に限るのか、「医師の診断」でよいのか。脳梗塞と脳出血を合わせて数えるのか、分けるのか。何年の追跡時点の値を使うのか。試験によっては「脳卒中、心筋梗塞、心血管死のいずれか」のような複合アウトカムしか報告していないこともあります。複合アウトカムは中身の内訳によって意味が変わるので、脳卒中だけの数が取り出せない試験をどう扱うかも、この段階で決めておきます。

\section{組み入れ基準と除外基準}
\label{sec:02-3}

PICO の 4 要素に、研究デザインなどの条件を加えて、どの研究を組み入れ、どの研究を除外するかの基準を作ります。降圧薬 A のレビューでは、たとえば次のように書きます。

\begin{description}
\item[研究デザイン] ランダム化比較試験に限ります。観察研究（薬剤を飲んでいた人と飲んでいなかった人を、ランダム化せずに比べた研究）は、両群の背景の違いが結果に混ざるため除外します。
\item[対象] 18 歳以上の高血圧の患者とします。高齢者に限った試験も、一般の成人の試験も組み入れます。
\item[介入と比較] 降圧薬 A とプラセボを比べた試験とします。他の降圧薬との比較試験は除外します。
\item[アウトカム] 脳卒中の発症を報告している試験とします。
\item[追跡期間] 1 年以上とします。短すぎる試験では脳卒中がほとんど起こらないためです。
\item[言語と出版形態] 言語は制限しません。学会抄録や試験登録の結果も探します。
\end{description}

言語を英語に限ると作業は楽になりますが、英語以外の言語で発表された試験が系統的に抜け落ちます。効果がはっきり出た試験ほど英語の国際誌に投稿されやすいという報告もあり、言語の制限は偏りの原因になりえます。翻訳の手間を見込んででも、言語は制限しないのが原則です。

基準は、迷ったときに 2 人の評価者が同じ判断に至るよう、具体的に書きます。「適切な追跡期間」ではなく「1 年以上」と書くのは、糖尿病型の判定を「空腹時血糖が高め」ではなく「126 mg/dL 以上」と数値で定めるのと同じ理由です。

\section{計画を先に書いて登録する}
\label{sec:02-4}

問いと基準が決まったら、データを集める前に、レビューの計画書（プロトコル）を書きます。計画書には、問い、組み入れ基準と除外基準、検索するデータベース、選別とデータ抽出の手順、バイアスのリスクの評価方法、統計解析の方法を書きます。プロトコルに書く項目は PRISMA-P という計画書の書式（Moher 2015）にまとめられているので、それに沿って書けば抜けが生じにくくなります。

計画書は、PROSPERO という国際的な登録のウェブサイトに登録するのが一般的です。計画を登録すると登録番号が付き、計画の内容と登録日が公開され、誰でも見られるようになります。

計画を先に書いて公開する理由は、結果を見てから方法を選ぶことを防ぐためです。メタ解析には、どの試験を入れるか、どの効果指標を使うか、どのモデルを使うかといった選択肢が数多くあります。結果を見ながら選ぶと、意識していなくても、都合のよい結果が出る選択に寄っていきます。検査の世界で、測定結果を見てから基準範囲の決め方を変えれば、どんな結果でも「正常」にできてしまうのと同じです。計画書を事前に登録しておけば、計画から外れた解析をしたときにそれが外から分かり、読者は結果をその分だけ割り引いて読むことができます。

データを集めてから初めて分かることもあるので、計画は途中で変えてもかまいません。変えたときは、いつ、何を、なぜ変えたかを記録し、論文に書きます。

\section{解析の計画で先に決めておくこと}
\label{sec:02-5}

統計解析の計画には、少なくとも次の項目を書きます。各項目の意味は第 5〜7 章で説明します。初めて読むときは、効果指標（\kkref{sec:05-2}）とモデル（\kkref{sec:05-4}）を先に決めることだけを押さえれば十分です。

\begin{kkcenter}
\kkfit{\begin{tabular}{ll}
\toprule
項目 & 降圧薬 A のレビューでの決め方 \\
\midrule
効果指標（第 5 章） & リスク比 \\
統合のモデル（第 5 章） & ランダム効果モデル（τ² は REML で推定） \\
異質性の評価（第 6 章） & I²、τ²、予測区間 \\
サブグループ解析（第 6 章） & 対象（高齢者と一般の成人） \\
感度分析（第 7 章） & バイアスのリスクが高い試験を除く \\
出版バイアス（第 7 章） & 試験数が 10 以上ならファネルプロットと Egger 検定 \\
確かさの評価（第 7 章） & GRADE \\
\bottomrule
\end{tabular}}
\kkcap{表 2.1　解析計画に書く項目と、降圧薬 A のレビューでの決め方}
\end{kkcenter}

サブグループ解析は、事前に決めたものだけを主な結果として扱います。データを見てから「この区分で分けたら差が出た」というサブグループは、偶然に見つかった差である可能性が高く、新しい仮説の候補にとどめます。表 2.1 で高齢者と一般の成人という区分を挙げたのは、高齢者では脳卒中の起こりやすさが高く、効果の大きさが違う可能性があると、データを見る前から考えられるからです。

\section{この章のまとめ}
\label{sec:02-6}

問いは PICO の 4 要素に分けて一文に書き、アウトカムは患者にとって意味のある出来事を選んで定義まで決めます。組み入れ基準は 2 人が同じ判断に至るよう具体的に書き、言語は制限しません。そのうえで、解析の方法まで含めた計画書を、データを見る前に書いて登録します。次の章では、この計画に沿って文献を探し、組み入れる試験を選びます。
